\documentclass[a4paper, 14pt]{extarticle}

\usepackage{geometry}
\geometry{a4paper,tmargin=2cm,bmargin=2cm,lmargin=3cm,rmargin=1cm}
\usepackage[T2A]{fontenc}
\usepackage[utf8]{inputenc}
\usepackage[english, main=russian]{babel}
\usepackage{textcomp}
\usepackage{indentfirst}
\usepackage{graphicx}
\graphicspath{{./images/}}
\usepackage{wrapfig}
\usepackage{float}
\usepackage[section]{placeins}
\linespread{1.3}
\sloppy
\clubpenalty=10000
\widowpenalty=10000
\usepackage{enumitem}
\usepackage{caption}
\usepackage{hyperref}
\hypersetup{unicode=true}
\DeclareCaptionLabelSeparator*{emdash}{~--- }
\captionsetup[figure]{labelsep=emdash,font=onehalfspacing,position=bottom}
\usepackage{tempora}
\usepackage{listings}
\lstset{
  basicstyle=\ttfamily\footnotesize,
  breakatwhitespace=false,
  breaklines=true,
  captionpos=t,
  inputencoding=utf8,
  frame=single,
  keepspaces=true,
  keywordstyle=\bf,
  numbers=left,
  numbersep=5pt,
  xleftmargin=25pt,
  xrightmargin=25pt,
  showspaces=false,
  showstringspaces=false,
  showtabs=false,
  stepnumber=1,
  tabsize=2,
  title=\lstname
}
\usepackage{amsthm,amsfonts,amsmath,amssymb,amscd}
\usepackage{mathtools}
\usepackage[perpage]{footmisc}

\begin{document}

\begin{titlepage}
\thispagestyle{empty}
\newpage
\vspace*{-60pt}
\hspace{-65pt}
\begin{minipage}{0.3\textwidth}
\hspace*{-20pt}\centering
\includegraphics[width=\textwidth]{emblem}
\end{minipage}
\begin{minipage}{0.67\textwidth}\small \textbf{
\vspace*{-0.7ex}
\hspace*{-6pt}\centerline{Министерство науки и высшего образования Российской Федерации}
\vspace*{-0.7ex}
\centerline{Федеральное государственное бюджетное образовательное учреждение }
\vspace*{-0.7ex}
\centerline{высшего образования}
\vspace*{-0.7ex}
\centerline{<<Московский государственный технический университет}
\vspace*{-0.7ex}
\centerline{имени Н.Э. Баумана}
\vspace*{-0.7ex}
\centerline{(национальный исследовательский университет)>>}
\vspace*{-0.7ex}
\centerline{(МГТУ им. Н.Э. Баумана)}}
\end{minipage}

\vspace{-25pt}
\hspace{-35pt}\rule{\textwidth}{2.3pt}
\vspace*{-20.3pt}
\hspace{-35pt}\rule{\textwidth}{0.4pt}

\vspace{1.5ex}
\hspace{-35pt} \noindent \small ФАКУЛЬТЕТ\hspace{80pt} <<Информатика и системы управления>>

\vspace*{-16pt}
\hspace{47pt}\rule{0.83\textwidth}{0.4pt}

\vspace{0.5ex}
\hspace{-35pt} \noindent \small КАФЕДРА\hspace{50pt} <<Теоретическая информатика и компьютерные технологии>>

\vspace*{-16pt}
\hspace{30pt}\rule{0.866\textwidth}{0.4pt}

\vspace{11em}

\begin{center}
\Large {\bf Лабораторная работа № 7} \\
\large {\bf по курсу <<Алгоритмы компьютерной графики>>} \\
\end{center}\normalsize

\vspace{8em}

\begin{flushright}
  {Студент группы ИУ9-41Б \\
  Хромов М.А. \\
  \\
  \vspace{2ex}
  Преподаватель Цалкович П.А.}
\end{flushright}

\bigskip
\vfill

\begin{center}
\textsl{Москва 2026}
\end{center}
\end{titlepage}

\renewcommand{\ttdefault}{pcr}
\newpage
\setcounter{page}{2}

\section{Задание}
Цель работы: изучение эффективных приемов организации приложений и оптимизации вызовов \texttt{OpenGL}.

В качестве основы выбрана лабораторная работа №6. Необходимо оптимизировать приложение с использованием не менее четырех различных методов, среди которых обязательно присутствуют дисплейные списки и массивы вершин. Также требуется провести измерения производительности и оформить результаты в табличном виде.

\section{Теория}
В лабораторной работе использованы следующие методы оптимизации:
\begin{enumerate}
  \item дисплейные списки;
  \item массивы вершин;
  \item индексированная отрисовка;
  \item отсечение обратных граней и уменьшение детализации модели во время анимации.
\end{enumerate}

\section{Реализация}
В программе предусмотрены пять режимов:
\begin{enumerate}
  \item baseline;
  \item display\_lists;
  \item display\_lists\_vertex\_arrays;
  \item display\_lists\_vertex\_arrays\_indexed;
  \item all\_optimizations.
\end{enumerate}

Замеры производительности выполняются автоматически и сохраняются в файл \texttt{lab7\_benchmark\_results.csv}.

\section{Результаты замеров}
После выполнения замеров таблица должна быть заполнена фактическими значениями.

\begin{table}[H]
\centering
\caption{Результаты замеров производительности}
\begin{tabular}{|l|c|c|c|}
\hline
Режим & Время, с & Средний FPS & Среднее время кадра, мс \\
\hline
baseline & --- & --- & --- \\
\hline
display\_lists & --- & --- & --- \\
\hline
display\_lists\_vertex\_arrays & --- & --- & --- \\
\hline
display\_lists\_vertex\_arrays\_indexed & --- & --- & --- \\
\hline
all\_optimizations & --- & --- & --- \\
\hline
\end{tabular}
\end{table}

\section{Заключение}
В ходе работы было выполнено исследование способов оптимизации приложения \texttt{OpenGL}, основанного на лабораторной работе №6. Для оценки применимости выбранных методов реализован автоматический механизм замеров производительности, позволяющий сравнить различные режимы рендеринга по среднему FPS и среднему времени кадра.

\end{document}
